\section{Introduction}
\label{sec:introduction}

Volatility plays a central role in modern financial markets. In its most basic interpretation, volatility measures the uncertainty of asset returns. In derivative pricing, however, volatility is more than a statistical input: it becomes a risk factor, a market expectation, and in many cases a directly tradable quantity. Products such as variance swaps, VIX futures, and VIX options allow investors to take positions on future realised or implied volatility without necessarily taking a direct directional position in the underlying equity index.

The development of volatility-linked derivatives has changed the way market participants understand and manage risk. A variance swap, for example, provides exposure to realised variance over a fixed time interval. VIX futures and VIX options provide exposure to market expectations of future volatility. These instruments are widely used for hedging, speculation, relative value trading, and volatility risk premium strategies. Their pricing therefore requires a framework that connects stochastic variance dynamics, risk-neutral valuation, numerical approximation, and empirical validation.

The main objective of this project is to study volatility as an asset class through variance swaps and VIX-linked derivatives. The project develops a quantitative framework that combines theoretical derivations with numerical pricing methods. The theoretical part is based on a Heston-type stochastic volatility model, where the instantaneous variance follows a Cox--Ingersoll--Ross square-root process. This model is useful as a benchmark because it captures mean reversion in variance while preserving a level of analytical tractability.

The project is organised around four completed components. The first component reviews the literature on volatility derivatives, variance risk premia, stochastic volatility models, and recent numerical methods. The second component develops the mathematical framework, including the variance swap strike, the model-implied VIX proxy, and the characteristic function of terminal CIR variance. The third component implements a Monte Carlo engine for Heston simulation, realised variance estimation, variance swap validation, and model-implied VIX call pricing. The fourth component implements a Fourier-cosine pricing engine for VIX derivatives using the characteristic function of terminal variance.

The numerical part of the project is deliberately built around two complementary methods. Monte Carlo simulation is flexible and intuitive. It allows stock and variance paths to be generated directly and can be used for path-dependent quantities such as realised variance. Its main limitation is that it contains sampling error and may require many paths for stable estimates. The Fourier-cosine method, in contrast, is deterministic once the model parameters and truncation interval are fixed. It can be computationally efficient when the characteristic function of the relevant state variable is available. The use of both approaches provides a natural cross-validation framework.

A key modelling distinction in this project is the difference between the official Cboe VIX index and the model-implied VIX proxy used in the numerical implementation. The official VIX is computed from a strip of out-of-the-money SPX option prices according to the Cboe methodology. The proxy used in this project is instead derived from the conditional expectation of future variance under the Heston model. Therefore, the project does not claim to reproduce the exchange methodology exactly. Rather, it studies a Heston-implied VIX-type quantity that is mathematically convenient for pricing and validation.

The research contribution of the current implementation is not a new stochastic volatility model, but a structured numerical framework that connects the theory and implementation of volatility-linked derivatives. The Monte Carlo engine provides a simulation-based estimate of realised variance and VIX option payoffs. The COS engine provides a characteristic-function-based benchmark. The comparison between the two engines helps verify whether the implementation is internally consistent.

The rest of the paper is structured as follows. Section~\ref{sec:literature_review} reviews the relevant literature and recent methodological developments. Section~\ref{sec:theory} presents the mathematical framework. Section~\ref{sec:monte_carlo} develops the Monte Carlo methodology. Section~\ref{sec:cos_pricing} presents the CIR-based COS pricing engine. The later sections discuss numerical validation, code architecture, limitations, and future extensions toward real data calibration and volatility risk premium backtesting.